\documentclass[11pt,a4paper]{article}

\usepackage[utf8]{inputenc}
\usepackage[T1]{fontenc}
\usepackage[french]{babel}
\usepackage[margin=2.3cm]{geometry}
\usepackage{enumitem}
\usepackage{booktabs}
\usepackage{tabularx}
\usepackage{array}
\usepackage{fancyhdr}
\usepackage{titlesec}
\usepackage{graphicx}
\usepackage{float}
\usepackage[most]{tcolorbox}
\usepackage{hyperref}
\hypersetup{colorlinks=true, linkcolor=black, urlcolor=black,
  pdftitle={Protocole experimental de collecte}}

\titleformat{\section}{\normalfont\Large\bfseries}{\thesection}{0.8em}{}
\titlespacing{\section}{0pt}{1.3em}{0.7em}
\titleformat{\subsection}{\normalfont\large\bfseries}{\thesubsection}{0.7em}{}
\titlespacing{\subsection}{0pt}{1.0em}{0.4em}

\pagestyle{fancy}\fancyhf{}
\fancyfoot[C]{\small Protocole expérimental de collecte --- rPPG / anémie}
\fancyfoot[R]{\small\thepage}
\renewcommand{\headrulewidth}{0pt}\renewcommand{\footrulewidth}{0.3pt}

\setlist[itemize]{leftmargin=*,itemsep=2pt,topsep=3pt}
\setlist[enumerate]{leftmargin=*,itemsep=3pt,topsep=3pt}

\newcolumntype{P}[1]{>{\raggedright\arraybackslash}p{#1}}

\newtcolorbox{notebox}[1]{colback=white,colframe=black,boxrule=0.6pt,arc=0pt,
  left=8pt,right=8pt,top=5pt,bottom=5pt,title={#1},coltitle=black,
  fonttitle=\bfseries\small,
  attach boxed title to top left={yshift=-2mm,xshift=2mm},
  boxed title style={colback=white,boxrule=0pt},breakable}
\newtcolorbox{warnbox}[1]{colback=white,colframe=black,boxrule=1pt,arc=0pt,
  left=8pt,right=8pt,top=5pt,bottom=5pt,title={#1},coltitle=black,
  fonttitle=\bfseries\small,
  attach boxed title to top left={yshift=-2mm,xshift=2mm},
  boxed title style={colback=white,boxrule=0pt},breakable}

\begin{document}

\begin{center}
{\LARGE\bfseries PROTOCOLE EXPÉRIMENTAL DE COLLECTE}\\[6pt]
{\large Screening cardiovasculaire et dépistage d'anémie\\ non invasifs par smartphone}\\[4pt]
{\normalsize Réseau de pharmacies --- Afrique subsaharienne}
\end{center}

\noindent\rule{\textwidth}{0.6pt}
\vspace{0.2cm}

% ================================================================
\section{Objectifs et motivations}
% ================================================================

L'objectif de cette collecte est de constituer un jeu de données de référence,
représentatif des conditions réelles de pharmacie africaine, afin d'entraîner et
d'évaluer des modèles de mesure de signes vitaux à partir d'une simple vidéo de
smartphone. Deux finalités complémentaires :

\begin{itemize}
\item \textbf{Mesure cardiovasculaire à distance (rPPG).} Les méthodes
d'estimation du rythme cardiaque par caméra ont été développées presque
exclusivement sur des populations à peau claire. Leur précision se dégrade
fortement sur les phototypes Fitzpatrick V--VI (la mélanine absorbe le signal
optique). Ce dataset donne une \textbf{priorité absolue aux phototypes V--VI}
pour permettre l'adaptation des modèles à cette population mal servie.
\item \textbf{Dépistage de l'anémie par imagerie.} L'anémie (taux d'hémoglobine
abaissé) se traduit par une \emph{pâleur} de la conjonctive (intérieur de la
paupière inférieure) et du lit de l'ongle. En photographiant ces zones et en
mesurant en parallèle le taux d'hémoglobine de référence, on constitue les
données nécessaires à un modèle estimant l'hémoglobine à partir d'une photo ---
un outil de dépistage non invasif particulièrement utile là où l'accès au
laboratoire est limité.
\end{itemize}

Chaque participant bénéficie en retour d'un \textbf{bilan de santé immédiat}
(section~\ref{sec:cloture}).

% ================================================================
\section{Mesures et données collectées}
% ================================================================

\begin{center}
\begin{tabularx}{\textwidth}{P{3.5cm} X P{4.2cm}}
\toprule
\textbf{Donnée} & \textbf{Description} & \textbf{Matériel} \\
\midrule
Vidéo du visage & Signal source rPPG (rythme cardiaque) & Smartphone, caméra avant, 60~fps \\
PPG / HR / SpO2 & Forme d'onde du pouls + fréquence cardiaque + saturation, référence & Oxymètre de pouls (USB) \\
Pression artérielle & Systolique / diastolique (mmHg) & Tensiomètre à brassard \\
\textbf{Hémoglobine} & Taux d'hémoglobine de référence (g/dL) & \textbf{Hémoglobinomètre} (prélèvement capillaire) \\
\textbf{Photos conjonctive} & Intérieur de la paupière inférieure (pâleur) & Smartphone, caméra \\
\textbf{Photos des doigts} & Lit de l'ongle / pulpe du doigt (pâleur) & Smartphone, caméra \\
Éclairage ambiant & Luminosité au niveau du visage (lux) & Luxmètre \\
\bottomrule
\end{tabularx}
\end{center}

\noindent Toutes les données sont enregistrées sous un \textbf{identifiant
anonyme} (ex.\ \texttt{CIV\_001}). Aucun nom ni coordonnée n'apparaît dans les
fichiers.

% ================================================================
\section{Matériel}
% ================================================================

\begin{itemize}
\item \textbf{Smartphones} (3 à 4 en rotation, modèles d'entrée de gamme
représentatifs du marché local) --- \textbf{caméra avant (selfie)}, comme
l'application de déploiement ; résolution et débit les plus élevés, compression
« intelligente » désactivée, \textbf{30~fps} (cadence de l'application),
\textbf{exposition et balance des blancs verrouillées}. L'\textbf{écran} peut servir
de \emph{lumière d'appoint} (mode blanc).
\item \textbf{Trépied / support smartphone} --- cadrage stable du visage (le patient
peut aussi tenir le téléphone à distance de bras, comme en usage réel).
\item \textbf{Anneau lumineux LED (optionnel)} --- utilisé \emph{uniquement} comme
\textbf{bras de comparaison} « lumière uniforme », \emph{pas} comme condition
principale : le déploiement étant une app en \textbf{lumière ambiante}, la majorité
des prises doit être en ambiance réelle variée (voir §\ref{sec:scenarios}).
\item \textbf{Oxymètre de pouls} (type Contec CMS50D+, USB) --- capture continue
PPG / HR / SpO2.
\item \textbf{Tensiomètre à brassard} (type Omron).
\item \textbf{Hémoglobinomètre} de point d'intervention (type HemoCue, mesure
capillaire au bout du doigt) --- fournit le taux d'hémoglobine de référence.
\item \textbf{Luxmètre} (capteur dédié ou application calibrée) --- seuil minimal
150~lux ; plage recommandée 200--800~lux pour les phototypes V--VI.
\item \textbf{Ordinateur portable + logiciel de collecte (interface graphique,
section~\ref{sec:gui})} --- pilote l'enregistrement synchronisé et la saisie des
métadonnées.
\end{itemize}

% ================================================================
\section{Population : priorité aux phototypes foncés}
% ================================================================

\begin{center}
\begin{tabularx}{\textwidth}{P{3.5cm} P{2cm} P{2.2cm} X}
\toprule
\textbf{Phototype} & \textbf{Part} & \textbf{Sujets} & \textbf{Rôle} \\
\midrule
Fitzpatrick I--III & 15\% & 10--15 & Comparaison avec la littérature \\
Fitzpatrick IV & 20\% & 15--20 & Population de transition \\
Fitzpatrick V & 30\% & 20--25 & Priorité haute \\
Fitzpatrick VI & 35\% & 25--30 & Priorité absolue \\
\bottomrule
\end{tabularx}
\end{center}

\noindent Le phototype est noté visuellement sur le terrain (échelle 1--6) pour
suivre les quotas, puis confirmé en post-traitement par un descripteur de teinte
de peau mesuré sur les pixels (ITA / vecteur de peau).

% ================================================================
\section{Scénarios d'enregistrement (format VitalVideos)}
\label{sec:scenarios}
% ================================================================

Le déploiement visé est une \textbf{application selfie} (caméra avant) utilisée en
\textbf{lumière ambiante non contrôlée}. La collecte doit donc \emph{ressembler au
déploiement} : le site principal est le \textbf{visage en selfie}, et la variable
déterminante --- surtout sur peaux foncées --- est l'\textbf{éclairage}. Plutôt que
de fixer une seule condition, on capture \textbf{plusieurs types de lumière} et on
\textbf{étiquette chaque vidéo dans le JSON}. Cela permet à la fois d'entraîner un
modèle robuste \emph{et} de \emph{stratifier} les analyses (p.\,ex. quantifier
l'apport réel d'une lumière uniforme sur les phototypes~V--VI).

\begin{center}
\begin{tabularx}{\textwidth}{P{2.6cm} P{5.6cm} X}
\toprule
\textbf{Champ (JSON)} & \textbf{Valeurs} & \textbf{Rôle} \\
\midrule
\texttt{site} & \texttt{visage} (principal) / \texttt{paume} (secondaire) & Visage = livrable ; paume = bras d'équité / recherche \\
\texttt{camera} & \texttt{avant} / \texttt{arriere} & \texttt{avant} = déploiement selfie \\
\texttt{lighting} & \texttt{ambiante\_bonne} / \texttt{\_moyenne} / \texttt{\_faible} / \texttt{naturelle\_fenetre} / \texttt{uniforme} & Ambiante = \textbf{majorité} (déploiement) ; uniforme = bras de comparaison \\
\texttt{screen\_fill} & \texttt{true} / \texttt{false} & Écran blanc en lumière d'appoint \\
\texttt{activity} & \texttt{repos} / \texttt{apres\_effort} & Élargit la plage de FC \\
\bottomrule
\end{tabularx}
\end{center}

\begin{notebox}{Répartition des conditions --- « domaine = déploiement »}
La \textbf{majorité} des prises doit être en \textbf{lumière ambiante réelle}
(\texttt{ambiante\_bonne/moyenne/faible}, \texttt{naturelle\_fenetre}), à divers
moments (matin/midi/soir, ensoleillé/nuageux) et en lumière \textbf{latérale}
(éviter le contre-jour) --- c'est ce que l'app verra. La lumière \textbf{uniforme}
(anneau) reste une \textbf{minorité}, comme bras de comparaison. \textbf{Inclure
délibérément des prises limites/sombres} (avec la référence oxymètre) : elles
servent à \textbf{calibrer l'abstention} (le seuil à partir duquel l'app doit dire
« mettez-vous à la lumière » plutôt que de rendre une valeur fausse). Chaque prise
est étiquetée par les champs ci-dessus dans le JSON.
\end{notebox}

% ================================================================
\section{Le logiciel de collecte (interface graphique)}
\label{sec:gui}
% ================================================================

La session est pilotée par une \textbf{interface graphique} dédiée, installée sur
l'ordinateur portable, qui automatise et fiabilise la collecte :

\begin{itemize}
\item \textbf{Enregistrement synchronisé} : elle déclenche la capture du signal
de l'oxymètre (PPG/HR/SpO2) et marque le début/fin de chaque enregistrement par
un \textbf{bip sonore} capté dans la piste audio de la vidéo du smartphone, ce
qui permet une synchronisation précise vidéo $\leftrightarrow$ signal de
référence en post-traitement.
\item \textbf{Affichage en temps réel} du tracé PPG, du HR et de la SpO2 ---
permet de vérifier la qualité du signal de référence avant de valider.
\item \textbf{Saisie du formulaire à la fin} de la session, écrivant les
\textbf{champs d'étiquetage} de chaque vidéo dans le JSON : phototype, âge, sexe,
\texttt{site} (visage/paume), \texttt{camera} (avant/arrière), \texttt{lighting}
(ambiante\_bonne/moyenne/faible/naturelle\_fenetre/uniforme), \texttt{screen\_fill},
\texttt{distance\_cm}, \texttt{activity} (repos/après-effort), \texttt{phone\_model},
FC de référence (médiane oxymètre), pression artérielle, taux d'hémoglobine, et
chemins des photos (conjonctive, doigts). Les champs \texttt{site}, \texttt{camera},
\texttt{lighting} et \texttt{activity} sont des \textbf{listes déroulantes} (valeurs
prédéfinies) pour garantir un étiquetage cohérent.
\item \textbf{Sortie structurée} : un dossier par participant, au format JSON
anonyme, regroupant vidéos, signal PPG, photos et métadonnées.
\end{itemize}

\begin{figure}[H]
\centering
\includegraphics[width=0.66\textwidth]{figures/gui_live.png}
\caption{Interface pendant l'enregistrement : tracé PPG en temps réel,
fréquence cardiaque et SpO2 de référence, numéro de patient anonyme, et bouton
de démarrage qui émet le bip de synchronisation.}
\end{figure}

\begin{figure}[H]
\centering
\includegraphics[width=0.66\textwidth]{figures/gui_form.png}
\caption{Formulaire rempli \emph{à la fin} de la session (phototype, scénario,
éclairage en lux, pression artérielle, etc.) --- saisi après l'enregistrement
pour ne pas perturber la mesure.}
\end{figure}

% ================================================================
\section{Procédure détaillée, étape par étape}
% ================================================================

\subsection{Étape 1 --- Accueil et consentement}
\begin{itemize}
\item Présentation orale du projet (en français et en langue locale si besoin) :
objectif, durée ($\sim$15~min), caractère anonyme et volontaire.
\item Lecture et signature du \textbf{formulaire de consentement éclairé}
(section~\ref{sec:consent}).
\item Attribution d'un identifiant anonyme (ex.\ \texttt{CIV\_001}) dans
l'interface de collecte.
\end{itemize}

\subsection{Étape 2 --- Installation}
\begin{itemize}
\item Sujet assis ; smartphone sur \textbf{trépied} à 50--70~cm, objectif au
niveau des yeux ; \textbf{anneau lumineux} en conditions contrôlées, ou éclairage
ambiant documenté au luxmètre.
\item Oxymètre placé au doigt du \textbf{bras~A} (sans brassard).
\item Brassard de tension sur le \textbf{bras~B}, \textbf{non gonflé} à ce stade.

\begin{warnbox}{Pourquoi séparer les deux bras}
Sur le même bras, le gonflage du brassard coupe momentanément la circulation et
fait perdre le signal de l'oxymètre. Les deux appareils sur des bras différents
permettent de capturer le PPG en continu pendant toute la session.
\end{warnbox}

\item Mesure du lux ambiant au niveau du visage (cible $\geq$150~lux).
\item Notation visuelle du phototype Fitzpatrick (1--6).
\item \textbf{Configurer la caméra} du smartphone (encadré ci-dessous).
\item Ouverture de la session dans le \textbf{logiciel de collecte}
(section~\ref{sec:gui}) sous l'identifiant anonyme.
\end{itemize}

\begin{notebox}{Réglages caméra (application OpenCamera) --- à faire une fois par appareil/session}
Le signal rPPG est une variation de couleur de $\sim$0{,}1~\%~: une compression
trop forte ou une exposition qui dérive le détruisent. \emph{Validation terrain~:
sur une même peau foncée (ITA $\approx -55$), passer aux réglages ci-dessous a fait
chuter le désaccord inter-méthodes de 15{,}9 à 0{,}3~bpm et l'erreur de 3{,}0 à
0{,}2~bpm.}
\begin{itemize}
\item \textbf{Résolution} : \textbf{1080p} (\emph{pas} 4K --- inutile pour le rPPG,
saturation mémoire, fichiers énormes).
\item \textbf{Débit (bitrate)} : « Custom » $\approx$ \textbf{50~Mbps}. Contrôle~:
un clip de 30~s doit peser \emph{plusieurs dizaines de Mo}~; $<$20~Mo = débit trop
faible.
\item \textbf{Cadence} : \textbf{60~fps} (ou 30 si lumière insuffisante).
\item \textbf{Anti-banding : 50~Hz} sous lumière artificielle (secteur africain
50~Hz) --- supprime le scintillement des lampes, qui peut se confondre avec un
pouls.
\item \textbf{Verrouiller l'exposition et la mise au point SUR LA PEAU} : taper sur
la \emph{joue} pour mesurer, puis activer le cadenas (AE/AF lock). Balance des
blancs sur une valeur \textbf{fixe} (pas « auto »).
\item \textbf{Compensation d'exposition (EV)} : \textbf{+0{,}3 à +1{,}0} pour les
phototypes V--VI si la peau est trop sombre (la remonter dans les tons moyens, sans
saturer).
\item \textbf{Désactiver} : HDR, stabilisation, réduction de bruit, « compression
intelligente / optimisation du stockage ».
\item \textbf{Anneau lumineux sans scintillement} (alimenté en continu / DC). Test~:
filmer l'anneau seul à 60~fps --- aucune bande ne doit défiler.
\end{itemize}
\end{notebox}

\subsection{Étape 3 --- Enregistrement 1 : état de repos (60--90~s)}
\begin{itemize}
\item Vérifier que le brassard reste \textbf{non gonflé}.
\item Dans le logiciel, bouton \textbf{«~DÉMARRER --- enregistrement 1 (repos)~»} :
il émet un \textbf{bip de début} et démarre la capture PPG/HR/SpO2 ; lancer en
même temps la vidéo du visage sur le smartphone.
\item Consigne selon le scénario en cours (section~\ref{sec:scenarios}) :
\texttt{faceonly} immobile et silencieux ; \texttt{facemovement} /
\texttt{handheld} mouvements naturels tolérés ; \texttt{slowbreath} /
\texttt{quickbreath} respiration guidée (lente ou rapide).
\item Bouton \textbf{«~FIN~»} : bip de fin. Vérifier immédiatement la qualité du
signal et de la vidéo ; en cas de problème, recommencer avant de continuer.
\end{itemize}

\subsection{Étape 4 --- Enregistrement 2 : pression artérielle (30--60~s)}
\begin{itemize}
\item Bouton \textbf{«~DÉMARRER --- enregistrement 2 (pression)~»} : bip de début ;
gonfler le brassard et lancer la mesure de pression \emph{pendant} cet
enregistrement. La vidéo et la capture PPG/HR/SpO2 continuent.
\begin{warnbox}{Mesurer la pression \emph{pendant} l'enregistrement}
Le gonflage du brassard peut provoquer une légère hausse transitoire du rythme.
Mesurer la pression pendant un enregistrement distinct (et non avant/après) garde
la cohérence entre la vidéo et l'état physiologique qu'elle représente.
\end{warnbox}
\item Bouton \textbf{«~FIN~»} : bip de fin.
\end{itemize}

\begin{notebox}{Les deux enregistrements dans un seul fichier}
Le logiciel regroupe les enregistrements~1 (repos) et~2 (pression) dans un
\textbf{même fichier JSON au format VitalVideos} (un scénario, deux
\emph{recordings}), la pression artérielle étant rattachée à l'enregistrement~2.
\end{notebox}

\subsection{Étape 5 --- Hémoglobine et photographies de dépistage}
\begin{itemize}
\item \textbf{Mesure du taux d'hémoglobine} à l'aide de l'hémoglobinomètre
(prélèvement capillaire au bout du doigt, microcuvette à usage unique).
\item \textbf{Photo de la conjonctive} : abaisser doucement la paupière
inférieure et photographier la conjonctive (zone interne rosée) sous éclairage
stable, cadrage rapproché et net.
\item \textbf{Photo des doigts} : photographier le lit de l'ongle et la pulpe des
doigts (main détendue, paume vers le haut), même éclairage.
\end{itemize}

\begin{notebox}{Conditions d'éclairage des photos}
La pâleur étant une information de couleur, l'éclairage des photos de conjonctive
et de doigts doit être le \emph{même} et documenté (idéalement sous l'anneau
lumineux). Une photo prise sous une lumière trop chaude ou trop faible fausse
l'évaluation de la pâleur.
\end{notebox}

\subsection{Étape 6 --- Saisie des mesures et enregistrement}
Après les deux enregistrements, le logiciel affiche le \textbf{formulaire de fin}.
Y reporter : phototype, position, mouvement, éclairage (lux), \textbf{pression
artérielle}, SpO2, \textbf{taux d'hémoglobine}, et les notes ; puis valider par
\textbf{«~ENREGISTRER~»}. Le logiciel écrit alors le fichier JSON complet (vidéos,
signal PPG des deux enregistrements, mesures, métadonnées) et \textbf{génère
automatiquement la fiche bilan} du participant.

\subsection{Étape 7 --- Clôture et remise du bilan de santé}
\label{sec:cloture}
\begin{itemize}
\item \textbf{Remise de la fiche bilan} (générée automatiquement par le logiciel,
imprimable), récapitulant : fréquence cardiaque (bpm), saturation en oxygène
(SpO2, \%), pression artérielle (mmHg), et taux d'hémoglobine (g/dL) avec
indication d'éventuelle anémie.
\item En cas de valeur anormale (hypertension, SpO2 basse, hémoglobine basse), le
participant est invité à consulter un professionnel de santé. \emph{Ce bilan est
informatif et ne constitue pas un diagnostic médical.}
\item Remerciement et compensation éventuelle selon les modalités convenues avec
la pharmacie.
\item Vérification finale que tous les fichiers (vidéos, signal PPG, photos,
métadonnées) sont enregistrés sous l'identifiant anonyme avant de passer au
participant suivant.
\end{itemize}

% ================================================================
\section{Formulaire de consentement éclairé}
\label{sec:consent}
% ================================================================

\noindent\textit{À lire et expliquer au participant ; une copie lui est remise
s'il le souhaite.}

\vspace{0.4cm}
\noindent\textbf{Objet de l'étude.}
Vous êtes invité(e) à participer à une étude visant à développer des outils de
mesure de signes vitaux (rythme cardiaque, saturation en oxygène, pression
artérielle) et de dépistage de l'anémie (taux d'hémoglobine) à partir de vidéos
et de photographies prises avec un smartphone. La participation dure environ
15~minutes.

\vspace{0.3cm}
\noindent\textbf{Ce qui sera enregistré.}
Une vidéo de votre visage, une courte vidéo de votre paume (optionnelle), des
photographies de votre conjonctive (intérieur de la paupière) et de vos doigts,
le signal de votre pouls et de votre oxygénation (oxymètre au doigt), votre
pression artérielle (brassard) et votre taux d'hémoglobine (petite goutte de sang
au bout du doigt).

\vspace{0.3cm}
\noindent\textbf{Utilisation et conservation des données.}
\begin{itemize}
\item Vos données seront utilisées pour \textbf{entraîner et évaluer des modèles}
d'intelligence artificielle de mesure de signes vitaux et de dépistage.
\item Elles seront \textbf{conservées de manière anonyme} : elles sont
enregistrées sous un identifiant codé, sans votre nom ni vos coordonnées.
\item Vos données anonymisées pourront être \textbf{partagées avec des
partenaires de recherche} à des fins scientifiques.
\end{itemize}

\vspace{0.3cm}
\noindent\textbf{Vos droits.}
\begin{itemize}
\item Votre participation est \textbf{entièrement volontaire}.
\item Vous avez le \textbf{droit de vous retirer de l'étude à tout moment}, sans
avoir à vous justifier et sans aucune conséquence.
\item Vous avez le \textbf{droit de poser toutes les questions} concernant le
traitement et l'utilisation de vos données, et d'obtenir une réponse claire.
\end{itemize}

\vspace{0.3cm}
\noindent\textbf{Bénéfice direct.}
À la fin de la session, un \textbf{bilan de santé} vous est remis (rythme
cardiaque, oxygénation, pression artérielle, taux d'hémoglobine). Ce bilan est
informatif et ne remplace pas un avis médical.

\vspace{0.6cm}
\noindent
\begin{tabularx}{\textwidth}{X X}
Identifiant anonyme : \underline{\hspace{3.5cm}} &
Date : \underline{\hspace{3.5cm}} \\[0.9cm]
Signature du participant : & Signature du collecteur : \\[1.0cm]
\underline{\hspace{6cm}} & \underline{\hspace{6cm}} \\
\end{tabularx}

\vspace{0.4cm}
\noindent\small
\textit{En signant, je confirme avoir compris l'objet de l'étude, l'utilisation
qui sera faite de mes données, et mes droits, et je consens librement à
participer.}

\end{document}
